\clearpage\begingroup\HistoricalHead{224}{L}{THEORETISCHE ASTRONOMIE. NACHLASS.}
\noindent{\small [S. 4]}
\begin{center}[4.]\end{center}
\noindent\hspace*{\parindent}Um aus diesen Datis die Bahn zu bestimmen, bedient man sich der Regula falsi, indem man einen so nahen Werth als möglich für den h[alben] Parameter, die Excentricität oder das Aphelium voraussetzt, daraus die [übrigen] Elemente der Bahn, und sodann die Zeit von der ersten zur zweiten Beobachtung ableitet, und mit der gegebenen Zwischenzeit vergleicht. Um aber bei der ersten Hypothese sich nicht zu sehr von der Wahrheit zu entfernen, bemerken wir, dass ziemlich genau $\surd p=\frac12(rr+r'r')\frac{v'-v}{U'-U}$ sei. Noch genauer erhält man $p$, wenn man aus dem ersten angenäherten Werthe die Grösse $\mathfrak{R}$ vermittelst der Formel
\[
\frac{1}{\mathfrak{R}}=
\frac{\frac12\left(\frac1r+\frac1{r'}\right)}{\cos\frac12(v'-v)}
-\frac{2\sin\frac14(v'-v)^2}{p\cos\frac12(v'-v)}
\]
\clearpage
\HistoricalHead{225}{R}{BESTIMMUNG DER BAHNEN DER HIMMELSKÖRPER.}\SourceLeaf{230}
\noindent oder
\[
\frac1{\mathfrak{R}}=\frac12\left(\frac1r+\frac1{r'}\right)
+\left(\frac12\left(\frac1r+\frac1{r'}\right)-\frac1p\right)
\frac{2\sin\frac14(v'-v)^2}{\cos\frac12(v'-v)}
\]
\noindent bestimmt u[nd] sich dann der Formel $\surd p=\frac16(rr+r'r'+4\mathfrak{R}\mathfrak{R})\frac{v'-v}{U'-U}$ bedient\textup{[\textsuperscript{*})]}, welche Operation man auch leicht mehreremal wiederholen kann. Auf diese Art bekommt man gewöhnlich, wenn die Zwischenzeit nicht zu gross ist, einen so genauen Werth von $p$, dass die folgende Rechnung mehr zur Bestätigung, als zur Verbesserung desselben dienen wird\textup{[\textsuperscript{**})]}.
\begingroup\renewcommand{\thefootnote}{}\footnotetext{[\textsuperscript{*}) Vgl. Werke VI, S. 161—162, ferner Theoria motus, art. 86, Werke VII, 1906, S. 109 und die Schlußworte des art. 84, a. a. O. S. 108.]}\endgroup
\begingroup\renewcommand{\thefootnote}{}\footnotetext{[\textsuperscript{**}) S. 8] Es sei $\dfrac{r'}r=\operatorname{tg}(45^{0}+\varphi)$, so dass
\[
\begin{gathered}
\frac1{rr}+\frac1{r'r'}=\frac2{rr'\cos2\varphi},\qquad
\frac1r+\frac1{r'}=\frac{2\cos\varphi}{\surd(rr'\cos2\varphi)};\\
\frac1r-\frac1{r'}=\frac{2\sin\varphi}{\surd(rr'\cos2\varphi)},\qquad v'-v=\delta,
\end{gathered}
\]
so ist
\[
\frac1{\mathfrak{R}}=
\frac{\cos\varphi\left(1-\dfrac{2\sin\frac14\delta^2}{p\cos\varphi}\surd(rr'\cos2\varphi)\right)}
{\cos\frac12\delta\surd(rr'\cos2\varphi)}
\]
und
\[
\begin{aligned}
\text{mot. m. }\odot\mathbin{.}\surd p&=\frac16\delta(rr+r'r'+4\mathfrak{R}\mathfrak{R})\\
&=\frac{\delta rr'}{\cos2\varphi}\left\{\frac13+\frac23
\left(\frac{\cos\frac12\delta\cos2\varphi}
{\cos\varphi\left(1-\dfrac{2\sin\frac14\delta^2}{p\cos\varphi}\surd(rr'\cos2\varphi)\right)}\right)^2\right\}
\end{aligned}
\]
oder näherungsweise [$m=\text{mot. m. }\odot$ gesetzt]
\[
p=\left(\frac{rr'\delta}{m}\right)^2\mathord{{}^{3}\!\surd}
\left\{\cos2\varphi^2\left(
\frac{\cos\frac12\delta}
{\cos\varphi\left(1-\dfrac{2\sin\frac14\delta^2}{p\cos\varphi}\surd(rr'\cos2\varphi)\right)}
\right)^8\right\}.
\]
}\endgroup
\begin{center}
\textls[100]{I. Bestimmung der Bahn, wenn $p$ als bekannt vorausgesetzt wird}\textup{[\textsuperscript{***})]}.
\end{center}
\begingroup\renewcommand{\thefootnote}{}\footnotetext{[\textsuperscript{***}) Vgl. Werke VI, S. 161 und Theoria motus, art. 79, Werke VII, 1906, S. 102.]}\endgroup
\[
\left\{\begin{aligned}
\cos v\mathbin{.}\left(1-\frac p{r'}\right)-\cos v'\mathbin{.}\left(1-\frac pr\right)
&=e\sin(v'-v)\mathbin{.}\sin\omega\\
\sin v'\mathbin{.}\left(1-\frac pr\right)-\sin v\mathbin{.}\left(1-\frac p{r'}\right)
&=e\sin(v'-v)\mathbin{.}\cos\omega
\end{aligned}\right\},
\]
\[
\begin{aligned}
e\sin(v'-v)=\surd\Bigl\{&\left(1-\frac pr\right)^2+\left(1-\frac p{r'}\right)^2\\
&-2\left(1-\frac pr\right)\left(1-\frac p{r'}\right)\cos(v'-v)\Bigr\},
\end{aligned}
\]
\[
\left\{\begin{aligned}
\left(1-\frac pr\right)\cos(v'-v)-\left(1-\frac p{r'}\right)
&=e\sin(v'-v)\mathbin{.}\sin(v-\omega)\\
\left(1-\frac pr\right)\sin(v'-v)
&=e\sin(v'-v)\mathbin{.}\cos(v-\omega)
\end{aligned}\right\},
\]
\HistoricalFoot{XI\textsubscript{1}.}{29}
\clearpage
\HistoricalHead{226}{L}{THEORETISCHE ASTRONOMIE. NACHLASS.}\SourceLeaf{231}
\[
\left\{\begin{aligned}
\left(1-\frac pr\right)-\left(1-\frac p{r'}\right)\cos(v'-v)
&=e\sin(v'-v)\mathbin{.}\sin(v'-\omega)\\
\left(1-\frac p{r'}\right)\sin(v'-v)
&=e\sin(v'-v)\mathbin{.}\cos(v'-\omega)
\end{aligned}\right\},
\]
\[
\frac{\dfrac p{r'}-\dfrac pr}{2\sin\frac12(v'-v)}=e\sin\bigl(\tfrac12(v'+v)-\omega\bigr),
\]
\[
\frac{2-\dfrac pr-\dfrac p{r'}}{2\cos\frac12(v'-v)}=e\cos\bigl(\tfrac12(v'+v)-\omega\bigr).
\]
\begin{center}
\textls[100]{II. Bestimmung der Bahn, wenn man $e$ [oder $\omega$] als bekannt voraussetzt}\textup{[\textsuperscript{*})]}.
\end{center}
\begingroup\renewcommand{\thefootnote}{}\footnotetext{[\textsuperscript{*}) Vgl. Werke VI, S. 160—161 und Theoria motus, art. 80—81, Werke VII, 1906, S. 103.]}\endgroup
\noindent\hspace*{\parindent}Allgemein. Man mache
\[
\begin{aligned}
e(r'\sin v'-r\sin v)&=k\sin\zeta,\\
e(r'\cos v'-r\cos v)&=k\cos\zeta,
\end{aligned}
\]
\noindent so ist $r'-r=k\cos(\omega-\zeta)$\textup{[\textsuperscript{**})]},
\begingroup\renewcommand{\thefootnote}{}\footnotetext{[\textsuperscript{**}) $\zeta$ hat hier eine andere Bedeutung, als in den folgenden Zeilen.]}\endgroup
\par\noindent folglich, wenn man macht
\[
\begin{aligned}
e(r'+r)\sin\tfrac12(v'-v)&=k\sin\zeta,\\
e(r'-r)\cos\tfrac12(v'-v)&=k\cos\zeta,
\end{aligned}
\]
\noindent so wird $r'-r=k\cos\bigl(\omega-\zeta-\tfrac12(v'+v)\bigr)$.
\par\noindent Also
\[
\operatorname{tang}\zeta=\frac{r'+r}{r'-r}\operatorname{tang}\tfrac12(v'-v),
\]
\[
\cos\bigl(\omega-\zeta-\tfrac12(v'+v)\bigr)=\frac{\cos\zeta}{e\cos\frac12(v'-v)}.
\]
\noindent{\small [S. 5]}\par
\noindent Eben diese Formel findet gleichfalls ihre Anwendung, wenn man $\omega$ als bekannt voraussetzt. Ferner ist sodann
\[
p=-\frac{2e\operatorname{tang}\zeta\cos\frac12(v'-v)\sin\bigl(\frac12(v'+v)-\omega\bigr)}
{\dfrac1r+\dfrac1{r'}}\text{[\textsuperscript{***})]}.
\]
\begingroup\renewcommand{\thefootnote}{}\footnotetext{[\textsuperscript{***})] Denomin. $=\dfrac{2\cos\psi}{\surd(rr'\cos2\psi)}$ ponendo $\dfrac r{r'}=\operatorname{tg}(45^{0}\pm\psi)$.}\endgroup
\clearpage
\HistoricalHead{227}{R}{BESTIMMUNG DER BAHNEN DER HIMMELSKÖRPER.}\SourceLeaf{232}
\begin{center}
\textls[100]{III. Bestimmung der mittlern Anomalie aus der wahren, und der Zeit.}
\end{center}
\[
\operatorname{tang}\tfrac12E=\operatorname{tang}\tfrac12M\mathbin{.}\operatorname{tang}(\tfrac12\varphi+45^{0}),
\]
\[
\begin{aligned}
\sin E&=\frac rb\sin M,\\
\cos E&=\frac ra\cos M-e=\frac{\cos M-e}{1-e\cos M},
\end{aligned}
\]
\[
\operatorname{cotang}\tfrac12(E-M)=\frac{\operatorname{cotang}\frac12\varphi}{\sin M}-\operatorname{cotang}M,
\]
\[
\sin(E-M)=\frac{er}{p}\sin M\mathbin{.}(1-\cos M\mathbin{.}\operatorname{tang}\tfrac12\varphi),
\]
\[
\sin\tfrac12(E-M)=\sin M\mathbin{.}\sin\tfrac12\varphi\surd\frac rp=\sin E\sin\tfrac12\varphi\surd\frac ar,
\]
\[
\sin\tfrac12(E+M)=\sin M\mathbin{.}\cos\tfrac12\varphi\surd\frac rp.
\]
\begin{center}\rule{.18\linewidth}{.35pt}\end{center}
\[
\sin(E'-E)=\frac{rr'}{ab}\sin(M'-M)-e\sin E'+e\sin E,
\]
\[
\sin\tfrac12(E'-E)^2=\frac{rr'}{bb}\sin\tfrac12(M'-M)^2,
\]
\[
\begin{aligned}
&\cos\tfrac12(E'-E)\cos\tfrac12(M'-M)
+\frac{\sin\frac12(E'-E)\sin\frac12(M'-M)}{\surd(1-ee)}\\
&\hspace{5em}=\frac{\surd\dfrac{r'}r+\surd\dfrac r{r'}}2,
\end{aligned}
\]
\[
\frac{\sin\frac12(E'+E)}{\sin\frac12(E'-E)}=\frac{\sin\frac12(M'+M)}{\sin\frac12(M'-M)}.
\]
\begin{center}\rule{.18\linewidth}{.35pt}\end{center}
\[
N=E+e\sin E;\qquad
\]
\[
N-M=2\operatorname{Arc}\sin\left\{\sin M\sin\tfrac12\varphi\surd\frac rp\right\}+\sin\varphi\sin M\mathbin{.}\frac rb.
\]
\noindent Zeit von der ersten zur zweiten Beobachtung $=\dfrac{a^{\frac32}(N'-N)}{\text{mot. diurn. }\odot\text{ med.}}$, wobei die mittl[ere] tägliche siderische Bewegung der Sonne $=3548^{\prime\prime}{,}19223$ und der Logar[ithmus] $=3{,}5500071$. — Stimmt diese berechnete Zeit $=T$ mit der beobachteten $t'-t$ nicht überein, so wiederholt man die Rechnung, indem man als neuen Werth des h[alben] Parameters $\dfrac{TT}{(t'-t)^2}\mathbin{.}p$ zum Grunde legt.
\clearpage\endgroup
